% Fragmento público generado desde propietarios íntegros.
% Residencia: 52.6.1.
% Fuente: ../incoming/radion_monografia_20260827/manuscrito/sections/07_incertidumbre_oscilador.tex:1-111
\noindent La forme symplectique coordonne action, incertitude et dynamique oscillatoire sans modifier la généalogie de l'état.\par\medskip

\subsubsection{Réalisation de l'état dans un espace de Hilbert}

\paragraph{Matrice de covariance minimale}

Une fois construite l'ellipse d'action, elle peut se réaliser dans un couple
d'opérateurs conjugués
\[
[\widehat Q,\widehat P]=i\hret.
\]
On définit
\begin{equation}
V_\eta=\frac{\hret}{2}
\begin{pmatrix}
e^\eta&0\\[2pt]0&e^{-\eta}
\end{pmatrix}.
\label{eq:covariance}
\end{equation}
Son déterminant est
\begin{equation}
\det V_\eta=\frac{\hret^2}{4}.
\label{eq:cov-det}
\end{equation}
Par conséquent,
\begin{equation}
(\Delta Q)^2=\frac{\hret}{2}e^\eta,
\qquad
(\Delta P)^2=\frac{\hret}{2}e^{-\eta},
\qquad
\Delta Q\,\Delta P=\frac{\hret}{2}.
\label{eq:uncertainty}
\end{equation}
L'inégalité de Robertson--Schrödinger est saturée.

L'annihilateur comprimé
\begin{equation}
\widehat a_\eta=
\frac{e^{-\eta/2}\widehat Q+i e^{\eta/2}\widehat P}
{\sqrt{2\hret}}
\label{eq:a-eta}
\end{equation}
satisfait \([\widehat a_\eta,\widehat a_\eta^\dagger]=1\). Son vide a pour
fonction d'onde
\begin{equation}
\psi_\eta(Q)=
(\pi\hret e^\eta)^{-1/4}
\exp\left(-\frac{Q^2}{2\hret e^\eta}\right).
\label{eq:gaussian}
\end{equation}

\paragraph{Ellipse comme niveau de la fonctionnelle de Mahalanobis}

L’ensemble
\begin{equation}
\mathbb E_{\mathrm{act}}
=\{x\in\R^2:x^{\mathsf T}V_\eta^{-1}x\le4\}
\label{eq:mahalanobis}
\end{equation}
a pour demi-axes
\begin{equation}
a_S=2\Delta Q,\qquad b_S=2\Delta P.
\label{eq:two-sigma}
\end{equation}
Son aire est
\[
4\pi\sqrt{\det V_\eta}=2\pi\hret=\hfull.
\]
L'ellipse HMT construite précédemment coïncide exactement avec le contour à deux
écarts de l'état gaussien minimal.

\begin{proposition}[Chaîne unique du quotient de forme]
Après avoir déclaré les interfaces spectrale, symplectique et de Hilbert,
\begin{equation}
\boxed{
e^{-\eta}
=\sqrt{\frac{A-C^*}{A+C^*}}
=\frac{b_1}{a_1}
=\frac{b_S}{a_S}
=\frac{\Delta P}{\Delta Q}.}
\label{eq:shape-chain-hilbert}
\end{equation}
\end{proposition}

L'ellipse d'action correspond au niveau classique \(H=\hret\omega\) et au
contour \(2\sigma\) du vide comprimé. Elle ne doit pas être confondue avec l'énergie
moyenne du vide, \(\hret\omega/2\). Cette différence d'un facteur deux est un
contrôle physique obligatoire.

\paragraph{Conditions d'invariance de l'ellipse comme objet opératoire}

Dans la formulation conventionnelle, une ellipse d'incertitude peut être considérée comme
un outil graphique pour représenter une matrice de covariance. Ici, l'ordre
causal est inversé : l'ellipse positive existe déjà comme publication de
\((A,C^*,\hret)\) ; l'espace de Hilbert la reconnaît au moyen de
\eqref{eq:covariance}. La réalité ontologique revendiquée ne signifie pas qu'une
particule classique parcourt physiquement un ovale dans le plan \((Q,P)\). Elle signifie
que l'état conserve une forme positive, une action totale et un quotient
d'orientation avant qu'une théorie probabiliste ne les interprète.

La fonction d'onde ne crée pas la forme ; elle la représente. L'incertitude ne crée pas
\(\hbar\) ; elle publie l'impossibilité de comprimer simultanément les deux
directions sans modifier le produit symplectique.

\begin{narrativebox}
La lecture probabiliste est une réalisation possible de l'ellipse, non sa cause.
L'objet premier est la conservation conjointe du produit et de l'orientation. La
probabilité apparaît lorsque cette géométrie est publiée comme covariance d'un
état sur Hilbert.
\end{narrativebox}

